% Replacement text for the revised body of the paper.
% Use these sections in place of the current Motivation/Research Gap,
% Technical Approach, Results, Discussion, and Conclusion text.

\subsection{Revised Motivation}

ECLIPSE was built to address the coordination challenge of lunar infrastructure design by connecting SysML-based system modeling, GIS-based site selection, urban planning, discrete-event simulation, CAD submission, and 3D visualization. Its original implementation proved that a lunar-base architecture could be selected, simulated, and visualized in a common environment. It also revealed a practical question that determines whether such a pipeline can support design iteration: when a stakeholder changes a mission input, does that change propagate through the data products and into the simulation and visualization used to evaluate the mission?

Three weaknesses motivated the work reported here. First, the visualization relied on a prepared playback structure even when upstream placement, routes, or simulation data changed. Second, CAD submission was limited to USD-family inputs and did not provide a consistent mechanism for preserving engineering metadata through conversion. Third, the DES was intentionally developed as an ISRU proof of concept, leaving much of its operating logic, equations, and assumptions embedded in code. The objective of this work was therefore not to replace every subsystem in ECLIPSE, but to strengthen these three links so that scenario definition, simulation, CAD geometry, and visualization are connected by explicit, traceable artifacts.

\subsection{Research Gap}

Existing digital-engineering pipelines often demonstrate interoperability by passing a selected scenario from one tool to the next. They are less often evaluated on whether a changed scenario can be regenerated, inspected, and compared without manually rebuilding downstream representations. ECLIPSE exhibited this distinction in three places. The simulation-to-visualization link required scenario-specific playback logic and did not provide a single formal package that Omniverse could load. CAD submission assumed USD-compatible geometry and did not distinguish the engineering metadata of the submitted artifact from metadata needed only to render it. Finally, the DES represented one ISRU logistics architecture whose roles, routes, thresholds, and power assumptions were partly embedded in Python and only partly exposed to the user.

The research gap is therefore an end-to-end one: a lunar Digital Engineering Ecosystem needs explicit contracts that preserve scenario identity, geometry authority, operational assumptions, and visual context as data move between tools. This work addresses that gap through scenario-scoped visualization packages, multi-format CAD ingestion with SysML metadata write-back, and a configurable ISRU mode complemented by an initial generic behavior framework.

\section{Technical Approach}

\subsection{A. Scenario-Scoped, Terrain-Aware Simulation-to-Omniverse Integration}

The original Omniverse playback was a useful demonstration but not a reusable integration contract. Module placement, rover timing, and route behavior were prepared for a known scenario, so a new mission configuration could require manual scene work. The revised workflow creates a manifest after the scenario has been simulated. The manifest references the scenario-specific scene USD, waypoint USD, terrain USD, system JSON, and DES result package; it also records actor definitions, movement data, and terrain-projection diagnostics. The package is stored under the scenario output directory, allowing several completed scenarios to coexist without sharing a single overwritten scene or log.

The Omniverse extension acts primarily as a package reader and playback client. It lists the available scenario manifests and loads the package chosen by the user. It creates or updates actors from the manifest, plays their state against the time-indexed DES log, and exposes route visibility together with overview, follow, and rover-chase camera modes. This removes the earlier requirement to encode the mission directly in the extension. Some compatibility logic remains for the established ISRU telemetry structure; the extension is therefore manifest-driven, rather than completely free of scenario-specific assumptions.

Spatial consistency is addressed through a common map frame and terrain sampling. Urban Planning provides the metric frame used to place modules and routes. The manifest builder aligns the prepared terrain asset to that frame, records the terrain transformation, and constructs a sampler from the terrain mesh. Static module footprints are sampled across their bottom surfaces; a local least-squares ground plane determines pitch and roll, and a final vertical offset keeps the transformed CAD footprint from penetrating the terrain. The CAD bounding box used for this operation includes internal transforms, avoiding case-by-case corrections for assets whose geometry was offset inside the source USD.

Routes are densified, projected to the mesh, and evaluated segment by segment. The extension can render projected routes with slope-dependent colors and waypoint markers. Rover yaw and pitch follow the local tangent of the projected path, while motion is parameterized by distance along the route rather than waypoint index. The DES uses scenario route lengths and selected route-level metrics, but it does not yet simulate traversal at every terrain waypoint. Similarly, the implemented workflow aligns the terrain asset supplied for a scenario; automatic retrieval and clipping of arbitrary site patches from a global lunar terrain database remains outside the current scope.

\subsection{B. Multi-Format CAD Ingestion and SysML Metadata Update}

The CAD submission objective was to make stakeholder geometry usable by the pipeline even when it was not delivered as USD. The implemented workflow accepts STEP/STP, STL, OBJ, and USD-family files. STEP/STP inputs are read through OpenCascade using pythonOCC and tessellated into meshes. STL and OBJ are read as meshes through trimesh. USD-family files are opened through the USD Python API. All supported formats are standardized into a USD asset for Omniverse and a GLB preview for the browser interface. The workflow is extensible through additional format readers, but it does not yet support every possible CAD format.

Conversion alone is insufficient because useful information can be absent from, changed by, or difficult to recover after conversion. The workflow therefore preserves two complementary metadata products. Source-CAD metadata describes the file as submitted: source name, file format, size, units when available, surface area, center of mass, and volume only when a watertight solid permits a reliable calculation. Visualization metadata is generated after the USD asset is available and records its resolved geometry and rendering state: generated USD and preview paths, bounding box, stage units, mesh count, materials when available, and selected up and front axes. The user may select up and front axes in the preview for converted formats that lack a reliable orientation convention.

SysML continues to govern the intended dimensions of each architectural module. The CAD bounding box is therefore scaled to fit the SysML target dimensions rather than overwriting them. The same length factor is applied to derived physical values: surface area scales with the square of the factor, volume with the cube, and center-of-mass coordinates with the transformed geometry. The scaled values, together with provenance such as the source file and format, are written into an auto-generated metadata block in the corresponding SysML part and mirrored in the associated JSON asset. The result is a traceable update of the authoritative model, not merely a visual conversion. The current interface can also reuse CAD files stored in the repository CAD library; full integration with an external curated metadata library remains future work.

\subsection{C. DES: From an Opaque ISRU Demonstration to a Configurable and Extensible Framework}

The original DES was appropriate as a proof of concept but unsuitable as a general mission-definition interface. It represented a predetermined ISRU chain: regolith excavation and hauling, ISRU processing, LOX accumulation, LOX transport, and selected power consumers. Several operating choices were embedded in Python, including logistics assumptions, thresholds, storage behavior, timing, and power events. The user could vary a small set of sliders, but could not inspect the full operational logic or connect a selected architecture to a different resource flow.

For the ISRU preset, those scenario-specific choices are now collected in a structured configuration passed from the interface to the DES workflow. The configuration records simulation duration, placed instances, rover counts, route assignments, payload capacities, route-derived distances, transport and processing parameters, storage settings, power generation, battery state, and module power-demand models. The UI exposes these choices through scenario architecture, route, equation, and DES configuration views. The executed configuration is included in the results package, so a completed run can be reproduced and saved with its named scenario.

The ISRU implementation itself was extended rather than discarded. It supports multiple regolith rovers and multiple LOX rovers, multiple ISRU plant instances, separate input storage at each plant, and dispatch of regolith vehicles toward the least-filled eligible plant. LOX delivery is assigned through the scenario route definition, and rover cycles include both outbound and empty return travel. The engine derives its haul lengths from the scenario routes generated upstream. It still represents each route in the DES by its effective length and selected route metrics, rather than by a full waypoint-by-waypoint vehicle dynamics model.

The power model was also made explicit. Solar generation power, battery capacity, initial charge, management interval, and module demand can be configured. A module demand can be represented as an average, a time profile, or a restricted equation. ISRU processing and LOX-storage energy are recorded as pending demands and passed to the power manager, addressing the earlier omission in which major ISRU energy use could be reported without contributing to the power-feasibility calculation. The model remains simplified: it is deterministic and does not yet represent detailed charging operations, failures, or a full electrical-network analysis.

To move beyond a single ISRU template, the workflow also introduces an initial generic scenario compiler. Step 4 serializes placed module instances, SysML interfaces, resource routes, rover assignments, module classes, equations, and power configuration. The compiler creates SimPy storage containers, processing processes, rover processes, and a power-accounting process. Classes are inferred from graph connectivity, ports, resource-flow direction, and relevant physical attributes, with a user override available when the inference is inappropriate. A Source supplies an outgoing resource, a Processor consumes an input event and produces configured outputs after a processing time, a Storage retains and relays a resource, a Consumer receives power or resources without transformation, a Transporter moves resources along an assigned route, and a PowerGenerator supplies the power model.

The equation builder is intentionally constrained. It uses a safe parser rather than arbitrary Python evaluation and accepts only supported mathematical operations and named variables. An equation affects the simulation only when its output maps to an implemented behavior, such as a resource output, stored-resource level, processing time, travel time, energy consumption, power demand, or power generation. The generic engine has been exercised with simple cargo-relay, ice-processing, and water-transfer scenarios, demonstrating that the architecture is no longer limited conceptually to LOX production. It should not yet be described as able to simulate every lunar mission. Complex multi-resource synchronization, richer storage constraints, stochastic operations, complete unit checking, and a broader library of operational behaviors remain open work.

\section{Results and Demonstrated Capabilities}

The Lunar Spaceport-1 ISRU case was used to demonstrate the integrated workflow. A named scenario contains its SysML and JSON state, CAD associations, Urban Planning outputs, DES configuration and results, and an Omniverse package. After the DES workflow completes, the scenario package contains independent scene, waypoint, manifest, and result files. The Omniverse extension can select that package rather than loading a single global playback. This provides reproducible scenario-level visualization and prevents a new run from silently replacing the visualization inputs of another saved scenario.

The terrain-aware playback demonstrates that module and rover placement no longer relies on arbitrary visual offsets. Module footprints are projected to the prepared terrain mesh, CAD geometry is transformed using its resolved bounding box and orientation metadata, and rover routes are lifted to the terrain. The route overlay exposes slope classes while the selected cameras and time-synchronized dashboard allow a user to inspect motion and quantitative state together. The dashboard is demonstrated most completely for the ISRU telemetry schema, including production, storage, power, and rover fields; actor-level dashboard fields are also defined through the manifest.

The CAD workflow was demonstrated with native USD-family assets and with STEP/STP, STL, and OBJ submissions. For each supported input, the pipeline produces a standardized USD asset, a browser preview, and metadata products. The submitted geometry is visually scaled to the SysML dimensions, while scaled physical metadata is written back to SysML. This preserves the distinction between authoritative architectural dimensions and derived information from a particular CAD submission.

The DES results demonstrate two levels of flexibility. In the ISRU preset, users can change the scenario architecture, route assignments, rover and plant multiplicity, transport and production assumptions, and power models without editing Python source. In the generic engine, simple non-ISRU examples use the same resource-route and class framework to simulate cargo relay, ice processing, and water transfer. The interface also saves named configurations, reloads them, and compares saved completed runs through selected resource histories and final measures of effectiveness. These comparison views are functional scenario-management tools, but their resource aggregation and presentation should be expanded as the generic resource vocabulary grows.

\section{Discussion}

For the automation question, the manifest establishes a concrete contract between upstream mission definition and Omniverse. A scenario run generates a package that the extension can select and play, reducing manual scene preparation. The process still depends on the GitHub-based workflow and a user pull in the local Omniverse repository; it is automated within that architecture, not a continuously deployed cloud service.

For the flexibility question, ECLIPSE has progressed substantially beyond the original opaque ISRU demonstration. The ISRU mode is now inspectable and configurable, and an initial generic compiler translates a restricted mission description into SimPy processes. This is stronger than a collection of sliders because it links module roles, equations, routes, and resource movement. It is nevertheless an extensible framework rather than a completed universal simulator. The generic behaviors and equation contracts must be expanded and validated before high-consequence lunar operations or arbitrary multi-resource architectures can be represented with confidence.

For coherence, the map-frame and terrain-projection work ties the prepared terrain mesh, module placement, CAD scale, and route playback to the same scenario package. The remaining limitation is terrain availability: fidelity depends on the resolution and preprocessing of the available mesh, and selecting a site in the UI does not yet trigger automatic extraction of a new terrain patch from a global terrain database. Similarly, the telemetry interface is synchronized with recorded simulation state during playback; it does not yet stream from a DES that is running concurrently.

\section{Conclusions and Future Work}

This work moves ECLIPSE from a manually coupled lunar-base demonstration toward a more traceable digital engineering workflow. Scenario-specific manifests now connect the simulation outputs, scene assets, terrain reference, routes, and actors consumed by Omniverse. CAD submission accepts several common source formats, preserves distinct source and visualization metadata, and writes scaled engineering metadata back to the SysML source of truth. The ISRU DES is now configurable and reproducible through the interface, while an initial role-based generic engine demonstrates how non-ISRU resource scenarios can be represented without creating a new dedicated Python controller for every example.

The next technical step is to mature the generic behavior language: richer multi-resource processing, capacity and queue models, unit-aware validation, stochastic events, and more detailed power and rover operations are needed before broad lunar mission coverage can be claimed. Future visualization work should connect UI site selection to automatic terrain-patch extraction and expand the dashboard into a more resource-agnostic mission-results interface. The repository CAD library can be connected to curated engineering and lunar metadata records. Finally, scenario saving and comparison provide a foundation for asynchronous collaboration, but real-time multi-user mission design will require persistent backend services, authorization, conflict handling, and synchronized shared state beyond the current GitHub Pages and GitHub Actions architecture.

% Suggested figure placement and truthful captions:
% IV.A: "Transition from a hardcoded playback to a scenario-scoped manifest
% package and Omniverse extension."
% IV.A: "Terrain-projected module placement and slope-classified rover route."
% IV.B: "Supported CAD conversion paths and dual metadata propagation to SysML."
% IV.C: "User-visible ISRU scenario definition and DES configuration."
% IV.C: "Role-based generic scenario compiler: a processor example."
% Results: "Time-synchronized DES telemetry during Omniverse playback."
% Results: "Saved scenario comparison of selected histories and final measures
% of effectiveness."
